% Inference: the interval that does not assume equal rates, the sampling term, coverage.

\section{Inference}
\label{sec:inference}

\subsection{An interval that does not assume the models are alike}

The audit draws labels by stratified sampling without replacement, and the pooled count in a
concordant cell sums labels from sub-populations whose true rates may differ. A textbook
binomial interval is exact for a single rate under sampling with replacement, and neither
condition holds here. Using one anyway would assume, inside the confidence statement, the
very homogeneity that Assumption~\ref{ass:h} is supposed to make explicit and testable.

We use a bound that survives both departures.

\begin{lemma}[Design-valid pooled bound]
\label{lem:kl}
Let $x_c$ be the total number of audited responses credited by the human reference in cell
$c$, summed over sub-populations, with $n_c$ labels in total and sampling weights
$w_{gc} = n_{gc}/n_c$. Let $\tilde p_c = \sum_g w_{gc}\, p_{gc}$. Then for every $t > 0$,
\[
\Pr\!\left(\frac{x_c}{n_c} \ge \tilde p_c + t\right) \le e^{-n_c \mathrm{KL}(\tilde p_c + t \,\|\, \tilde p_c)},
\qquad
\Pr\!\left(\frac{x_c}{n_c} \le \tilde p_c - t\right) \le e^{-n_c \mathrm{KL}(\tilde p_c - t \,\|\, \tilde p_c)} .
\]
\end{lemma}

The proof, in Appendix~\ref{app:proofs}, combines two standard facts. Sampling without
replacement is dominated by sampling with replacement in the sense required for moment
generating functions, which removes the first difficulty, and the moment generating function
of a sum of independent Bernoulli variables with different means is bounded by that of a
single Bernoulli at the weighted mean, which removes the second. Inverting the bound gives an
interval $[L_c, U_c]$ with coverage at least the nominal level, valid in finite samples, for
heterogeneous rates, and without replacement.

Clopper--Pearson intervals are tighter, and we report them alongside as a sensitivity.
They are exact only under assumptions this design does not satisfy, so the headline numbers
use Lemma~\ref{lem:kl}.

\subsection{Item sampling}

The benchmark's \NumItems{} items are themselves a sample of the questions one might have
asked. A claim about which model is better in general, rather than on these items, carries a
second source of uncertainty that the audit cannot reduce. We bound it by a paired argument.
For a pair of models, let $D_i$ be the difference of their human-reference verdicts on item
$i$. Then $\mathrm{Var}(D) \le \Pr(D \neq 0)$, and $\Pr(D \neq 0)$ is bounded above by the
observed rate at which a verifier's verdicts differ between the two models, plus each model's
maximum possible disagreement with that verifier, both of which the identified range already
supplies. The resulting half-width is paired, so it exploits the fact that both models answer
the same items, and it is conservative, because it uses upper bounds rather than estimates.

This term is asymptotic. It rests on a central limit approximation, and we say so rather than
presenting the combined interval as exact. Appendix~\ref{app:inference} reports a
distribution-free alternative, which is valid in finite samples and noticeably wider.

\subsection{What is certified, and at what level}

Two sources of error are combined by a union bound, and the two are distinct. The audit
interval concerns the unknown correctness rates given these responses; the item-sampling term
concerns generalisation beyond these items. Nothing is counted twice.

Families of claims are controlled explicitly. A statement about one ordering is made at the
per-ordering level. A statement about the whole board is a statement about \NumOrderings{}
orderings at once, so the sampling term is widened by a Bonferroni factor over the family, and
the pooled concordant intervals are shared by every pair, which makes the audit part
simultaneous automatically. Rank intervals follow: if every certified ordering is correct,
model $m$'s rank is at least one plus the number of models certified above it and at most the
total less the number certified below it. We report rank intervals as the family-level object,
because a reader looking at a leaderboard is reading all the orderings at once, not one.

\begin{proposition}[Coverage]
\label{prop:coverage}
Under Assumption~\ref{ass:h}, all audit-only claims hold simultaneously with probability at
least $1 - \alpha_{\mathrm{audit}}$, exactly in finite samples. Claims that also generalise
beyond the benchmark's items hold with probability at least
$1 - \alpha_{\mathrm{audit}} - \alpha_{\mathrm{sample}}$, asymptotically in the number of
items.
\end{proposition}

The two halves of that statement are deliberately different in kind, and we do not smooth over
the difference. The first is a finite-sample guarantee that follows from Lemma~\ref{lem:kl}.
The second inherits the central limit approximation of the sampling term.

\subsection{Simulation}

Simulation cannot establish a coverage theorem, and we do not use it as one. It checks
whether the guarantee is attained, and how it degrades when the assumption fails. We simulate
at the full scale of the board, \NumModels{} models and \NumOrderings{} orderings, with
\SimReplicates{} replicates per scenario, drawing audits under the real design and the real
cell structure.

Where Assumption~\ref{ass:h} holds, no certified claim was wrong in any scenario, so the
procedure attains its stated level and is conservative rather than tight. Where the truth
carries more heterogeneity than the analysis allows for, the family-wise error rises to
\SimFwerMisspecified{}, against a nominal \SimWorstFwerUnderH{} observed under a correct
specification. A world in which every model shares the same underlying rate, so that the
per-model fractions differ only by chance, pushed the family-wise error to
\SimFwerPropensityKTwo{} when analysed at $\varepsilon = 0$, and a tolerance of \EpsPrereg{}
restored it. That last scenario is the quantitative reason the paper does not report
$\varepsilon = 0$ as a result.
